\subsection{傅里叶变换的连续性}

命题 8. --- 设 E 是局部凸空间，\(\mu\) 是 E 上的投影测度，\(\Phi\) 是 \(\mu\) 的傅里叶变换。有不等式

\[
\left| \Phi (x ^ {\prime}) \right| \leqslant \Phi (0)\tag{53}
\]

\[
\left| \Phi (x ^ {\prime}) - \Phi (y ^ {\prime}) \right| ^ {2} \leqslant 2 \Phi (0) \big (\Phi (0) - \mathcal {R} \Phi (x ^ {\prime} - y ^ {\prime}) \big)\tag{54}
\]

对于 \(x', y'\) 中的 \(\mathbf{E}'\)。

第~3号的公式 (5) 允许归结为 E 有限维且 \(\mu\) 是测度的情形。于是

\[
\left| \Phi (x ^ {\prime}) \right| = \left| \int_ {\mathrm{E}} e ^ {i \langle x, x ^ {\prime} \rangle} d \mu (x) \right| \leqslant \int_ {\mathrm{E}} | e ^ {i \langle x, x ^ {\prime} \rangle} | d \mu (x) = \int_ {\mathrm{E}} d \mu (x) = \Phi (0),
\]

因此得到 (53)。此外，如果 \(a\) 和 \(b\) 是实数，则

\[
\left| e ^ {i a} - e ^ {i b} \right| ^ {2} = \left| e ^ {i b} \right| ^ {2} \left| e ^ {i (a - b)} - 1 \right| ^ {2} = \left(e ^ {i (a - b)} - 1\right) \left(e ^ {- i (a - b)} - 1\right) = 2 - 2 \cos (a - b);
\]

由柯西--施瓦茨不等式，于是有

\[
\begin{array}{r l} \big | \Phi (x ^ {\prime}) - \Phi (y ^ {\prime}) \big | ^ {2} & = \Big | \int_ {\mathrm{E}} \big (e ^ {i \langle x, x ^ {\prime} \rangle} - e ^ {i \langle x, y ^ {\prime} \rangle} \big) d \mu (x) \Big | ^ {2} \\ & \leqslant \int_ {\mathrm{E}} \big | e ^ {i \langle x, x ^ {\prime} \rangle} - e ^ {i \langle x, y ^ {\prime} \rangle} \big | ^ {2} d \mu (x) \int_ {\mathrm{E}} 1 ^ {2} d \mu (x) \\ & = \int_ {\mathrm{E}} (2 - 2 \cos \langle x, x ^ {\prime} - y ^ {\prime} \rangle) d \mu (x) \cdot \Phi (0) \\ & = 2 \Phi (0) \big (\Phi (0) - \mathcal {R} \Phi (x ^ {\prime} - y ^ {\prime}) \big), \end{array}
\]

因此得到 (54)。

推论. --- 给 E' 赋予一个与其向量空间结构相容的拓扑。Φ 连续的充要条件是其实部 \(\Phi\)\(\mathcal{R}\Phi\) 在原点连续；在此情况下，Φ 一致连续。\(\Phi\)

这由不等式 (54) 得出。

设 F 是局部凸空间。给 F 的对偶 \(F'\) 赋予一个与 F 和 \(F'\) 之间的对偶关系相容的拓扑，并将 F 认作 \(F'\) 的对偶。因此，\(\mu\) 上有界测度 \(F'\) 的傅里叶变换是 F 上的函数 \(F\mu\)，由下式定义

\[
(\mathcal {F} \mu) (x) = \int_ {\mathrm{F} ^ {\prime}} e ^ {i \langle x, x ^ {\prime} \rangle} d \mu (x ^ {\prime}).
\]

命题 9. --- 如果 F 是桶状空间，则 \(F'\) 上每个有界测度的傅里叶变换都是 F 上的一致连续函数。

设 \(\mu\) 是 \(\mathbf{F}'\) 上的有界测度，\(\Phi\) 是其傅里叶变换。令 \(\varepsilon > 0\)。存在  的紧子集 \(\mathbf{K}\)，使得 \(\mathbf{F}'\)\(\mu(\mathbf{F}' - \mathbf{K}) \leqslant \varepsilon\)。现在，\(\mathbf{K}\) 对弱拓扑 \(\sigma(\mathbf{F}', \mathbf{F})\) 是紧的，因而由于 \(\mathbf{F}\) 是桶状空间而等度连续（TVS, III, §4, 第~2号，定理 1）。因此存在  中 0 的对称邻域 \(\mathbf{U}\)，其极集 \(0\)\(\mathbf{F}\)\(\mathbf{U}^{\circ}\) 包含 \(\mathbf{K}\)。令 \(x\)；则\(\varepsilon \mathrm{U}\)

\[
\Phi (0) - \mathcal {R} \Phi (x) = \int_ {\mathrm{F} ^ {\prime}} \left(1 - \cos \langle x, x ^ {\prime} \rangle\right) d \mu (x ^ {\prime}).
\]

现在，\(0 \leqslant 1 - \cos \langle x, x' \rangle \leqslant 2\) 对每个 \(x' \in \mathbf{F}' - \mathbf{K}\)，并且

\[
1 - \cos \langle x, x ^ {\prime} \rangle \leqslant \frac {1}{2} \langle x, x ^ {\prime} \rangle^ {2} \leqslant \frac {\varepsilon^ {2}}{2}
\]

对于 \(x^{\prime}\in \mathbf{K}\subset \mathbf{U}^{\circ}\)；所以

\[
0 \leqslant \Phi (0) - \mathcal {R} \Phi (x) \leqslant 2 \mu (\mathrm{F} ^ {\prime} - \mathrm{K}) + \frac {\varepsilon^ {2}}{2} \mu (\mathrm{K}) \leqslant 2 \varepsilon + \frac {\varepsilon^ {2}}{2} \mu (\mathrm{F} ^ {\prime}).
\]

这个不等式右端随 \(\varepsilon\) 趋于 0；因此 \(\mathcal{R}\Phi\) 在 0 连续，命题由命题 8 的推论得证。

